% ===============================================================
% ===============================================================
% BOLETÍN SMCCA (PLANTILLA/TEMPLATE)
% ===============================================================
% ES: Selecciona el idioma descomentando UNA de las dos líneas.
% EN: Select the language by uncommenting ONE of the two lines.

\documentclass[spanish]{smcca}  % ES: Artículo en español
% \documentclass[english]{smcca} % EN: Article in English

% ===============================================================
% ES: PAQUETES ADICIONALES (OPCIONAL)
% EN: ADDITIONAL PACKAGES (OPTIONAL)
% ===============================================================
% ES: Aquí puedes añadir paquetes extra solo si son necesarios.
% EN: Add extra packages here only if they are required.
%
% \usepackage{...}


% ===============================================================
% ES: DATOS DEL ARTÍCULO
% EN: ARTICLE INFORMATION
% ===============================================================

% ES: Escribe el título SIN forzar negritas (\textbf{}).
% EN: Write the title WITHOUT forcing bold (\textbf{}).
\title{Title / Título}

% ES/EN: Autores y afiliaciones.
% ES: Usa índices [1],[2],... para evitar repetir instituciones.
% EN: Use indices [1],[2],... to avoid repeating institutions.
%
% ES: Correos en la primera página con \thanks{} (uno por autor o agrupados).
% EN: Emails as first-page footnotes using \thanks{} (one per author or grouped).

\author[1,2]{First Author / Primer Autor}
\author[1]{Second Author / Segundo Autor\thanks{\texttt{a1@uni.mx}}}
\author[2]{Third Author / Tercer Autor}

\affil[1]{Institution 1, City, Country. / Institución 1, Ciudad, País.}
\affil[2]{Institution 2, City, Country. / Institución 2, Ciudad, País.}

% ES: Déjalo vacío.
% EN: Leave empty.
\date{}

\setcounter{page}{1}
\begin{document}
\maketitle
	
	% ===============================================================
	% ES: RESUMEN / EN: ABSTRACT
	% ===============================================================
	% ES: Máximo ~10 líneas. Incluye el objetivo, el método, los resultados y las conclusiones.
	% EN: Max ~10 lines. Include objective, method, main results, and conclusions.
	
	\begin{abstract}
		Abstract / Resumen.
	\end{abstract}
	
	\vspace{1em}
	
	% ES/EN: Palabras clave / Keywords
	% ES: 4–6 palabras clave separadas por punto y coma.
	% EN: 4–6 keywords separated by semicolons.
	\noindent\textbf{Palabras clave / Keywords:} Keyword 1; Keyword 2; Keyword 3; Keyword 4.
	
	\vspace{10pt}

	% ===============================================================
	% ES: CUERPO DEL ARTÍCULO / EN: MAIN TEXT
	% ===============================================================
	
	% ES/EN: Cambia los títulos de las secciones según el idioma y las necesidades del artículo.
	% ES: Ej. "Introducción", "Metodología", "Resultados", "Conclusiones".
	% EN: E.g. "Introduction", "Methodology", "Results", "Conclusions".
	
	\section{Introducción / Introduction}
		Text...
	
	\section{Metodología / Methodology}
		Text...
	
	\section{Resultados / Results}
		Text...
	
	\section{Conclusiones / Conclusions}
		Text...
	
	% ===============================================================
	% ES: BIBLIOGRAFÍA / EN: REFERENCES
	% ===============================================================
	% ES: Usa un archivo .bib (p. ej., bibliografia.bib) y cita con \cite{...}.
	% EN: Use a .bib file (e.g., bibliografia.bib) and cite with \cite{...}.
	
	% Se agrega el comando \nocite{*} para citar todas las referencias y evitar
	% el error "no \citation commands"; esto puede removerse en la versión final
	\nocite{*}
	
	\bibliographystyle{abbrv}
	\bibliography{bibliografia} 
\end{document}
